\begin{pastexam}[label={px:nchu-115-is-04}]{中興資工所乙組 115 學年度 第 4 題}
（單選；每題 2 分，答錯不倒扣）

Which operation is most difficult in a singly linked list?

A.~Insertion at head \qquad B.~Insertion after given node\\
C.~Deletion of tail node \qquad D.~Traversal
\end{pastexam}

\begin{pastsoln}
答案是 C。

設串列有 $n$ 個節點，有 $\mathit{head}$ 指標（也可以有 $\mathit{tail}$ 指標）。
\begin{itemize}
\item A：新節點的 $\mathit{next}$ 指向原來的頭，再把 $\mathit{head}$ 改成新節點，$O(1)$。
\item B：已經拿到節點 $x$ 的指標，新節點的 $\mathit{next}$ 設成 $x.\mathit{next}$，再把
$x.\mathit{next}$ 指向新節點，$O(1)$（資結書第~5.1~節）。
\item C：刪尾端要把「倒數第二個節點」的 $\mathit{next}$ 改成 \NIL。 單向串列只能往後走，就算有
$\mathit{tail}$ 指標，也不知道它前一個是誰，只能從頭走 $n - 1$ 步找到倒數第二個，$O(n)$。 這是
四個操作裡唯一一個「本來可以很快、卻因為單向而變慢」的，所以是最難的。 雙向串列有
$\mathit{prev}$ 指標，刪尾端就是 $O(1)$（資結書第~5.2~節）。
\item D：走訪本來就要看每個節點一次，$\Theta(n)$ 是它的本分，單向串列做起來和陣列一樣直接，
沒有額外的困難。
\end{itemize}
\end{pastsoln}
